\section{Conclusiones}
En este trabajo, se ha abordado el complejo problema de la simulación de una maniobra de asistencia gravitatoria, 
un pilar fundamental de la astrodinámica moderna; para ello, se ha desarrollado desde cero un 
simulador del Sistema Solar en C, utilizando un integrador Runge-Kutta de cuarto orden con paso de tiempo adaptativo, 
lo que ha permitido obtener un equilibrio muy bueno, creo yo, entre precisión numérica y coste computacional.

El objetivo principal, que era encontrar los ángulos de lanzamiento óptimos para que una sonda, 
con condiciones iniciales inspiradas en la misión Solar Orbiter, maximizara la transferencia de energía con 
Venus, se ha logrado con éxito; mediante un sistema de barrido de parámetros orquestado por Python y ejecutado por C, 
se ha podido explorar sistemáticamente el espacio de soluciones, identificando no una, sino dos familias de 
trayectorias óptimas.

Este proceso de optimización, que se ha llevado a cabo en varias fases, primero con un barrido grueso 
para mapear todo el espacio de posibilidades, y luego con sucesivos barridos finos de resolución creciente, 
ha demostrado ser una estrategia muy potente; ha permitido, de hecho, localizar con una precisión de hasta 0.005 
grados los ángulos de guiñada y cabeceo que minimizan la distancia de aproximación al planeta, que, como se demostró 
en la introducción, es la clave para una asistencia gravitatoria efectiva.

Los resultados, además, son físicamente muy interesantes; se ha encontrado una trayectoria que, efectivamente, 
pierde energía al interactuar con Venus, cumpliendo el objetivo de "frenar" para acercarse al Sol
($\Delta E \approx -2.9 \cdot 10^{10}$ J), pero también 
se ha identificado otra que, por el contrario, gana energía, alejándose (con un $\Delta E \approx +1.7 \cdot 10^{11}$ J). 
El análisis de las gráficas de energía de la nave confirma, de forma visual e inequívoca, estos saltos energéticos 
justo en el momento del sobrevuelo; mientras que la conservación de la energía total del sistema, con un error 
relativo del orden de $10^{-14}$, valida la robustez y estabilidad del integrador numérico implementado.

No obstante, este trabajo, como todo modelo, posee sus limitaciones inherentes; la ausencia de propulsión activa 
en la nave, la simplificación de la atmósfera de Venus a un umbral de colisión, o la no consideración de perturbaciones 
gravitatorias de cuerpos menores, son factores que, sin duda, podrían refinarse en futuras investigaciones; 
abriendo, así, nuevas vías para explorar optimizaciones multi-objetivo o, incluso, el diseño de complejas 
cadenas de asistencias gravitatorias, que es un campo, por cierto, apasionante y que pensaba explorar con este 
trabajo, pero no me ha sido posible.

En definitiva, se ha construido una herramienta de simulación y análisis muy completa, que no solo ha 
permitido recrear un fenómeno físico complejo con gran fidelidad, sino que también ha servido para aplicar 
y visualizar conceptos fundamentales de la mecánica clásica, como son los problemas de dispersión, la importancia de las condiciones iniciales para evitar trayectorias catastróficas, y la 
conservación de la energía y el momento angular, todo ello, además, con un resultado que, visualmente, es 
bastante espectacular.